\documentclass{notaaula}

        \titulonota{Fluidos em repouso e movimento}
        \creditos{Turma 1B · Física · Novembro/2026 · Prof. Josué Cavalcante}

        \begin{document}
        \cabecalho

        \begin{sobre}
        Noções básicas de hidrostática e hidrodinâmica, com situações de pressão, empuxo e escoamento. Habilidades em foco: EM13CNT101, EM13CNT306 e EM13CNT307.
        \end{sobre}

        \section{Densidade e pressão}

\begin{copiar}{Grandezas básicas}
A densidade relaciona massa e volume; a pressão relaciona força normal e área:
\[
  \rho=\frac{m}{V},
  \qquad
  p=\frac{F}{A}.
\]
No SI, densidade é medida em $\mathrm{kg/m^3}$ e pressão em pascal:
$1\un{Pa}=1\un{N/m^2}$.
\atencao{A mesma força produz maior pressão quando atua em uma área menor.}
\diaadia{Agulhas perfuram com facilidade porque concentram a força em uma área muito pequena.}
\end{copiar}

\begin{exemplo}
Uma força de $600\un{N}$ atua sobre área de $\dec{0,020}\un{m^2}$.
\[
  p=\frac{600}{\dec{0,020}}=
  \resultado{3\times10^4\un{Pa}}.
\]
\end{exemplo}

\section{Pressão em líquidos}

\begin{copiar}{Lei de Stevin}
Em um líquido homogêneo em repouso, a pressão aumenta com a profundidade:
\[
  p=p_0+\rho gh,
  \qquad
  \Delta p=\rho g\Delta h.
\]
Pontos no mesmo líquido e na mesma profundidade têm a mesma pressão. A forma do recipiente
não altera $\rho gh$.
\simbolos{$p_0$ é a pressão na superfície; $h$ é a profundidade.}
\end{copiar}

\begin{copiar}{Pascal e Arquimedes}
Em um fluido confinado, uma variação de pressão é transmitida a todos os pontos:
\[
  \frac{F_1}{A_1}=\frac{F_2}{A_2}.
\]
Um corpo imerso recebe empuxo para cima igual ao peso do fluido deslocado:
\[
  E=\rho_{\mathrm{fluido}}gV_{\mathrm{deslocado}}.
\]
O corpo flutua em equilíbrio quando $E=P$.
\end{copiar}

\begin{exemplo}[Prensa hidráulica]
O êmbolo menor tem área $5\un{cm^2}$ e recebe $100\un{N}$. O maior tem $200\un{cm^2}$.
\[
  F_2=F_1\frac{A_2}{A_1}=100\frac{200}{5}
  =\resultado{4000\un{N}}.
\]
A prensa multiplica a força, mas o êmbolo menor percorre distância maior: energia não é criada.
\end{exemplo}

\section{Fluidos em movimento}

\begin{copiar}{Vazão e continuidade}
A vazão volumétrica é o volume que atravessa uma seção por unidade de tempo:
\[
  Q=\frac{\Delta V}{\Delta t}=Av.
\]
Para escoamento permanente de líquido incompressível,
\[
  A_1v_1=A_2v_2.
\]
Assim, onde o tubo fica mais estreito, a rapidez do líquido aumenta. Em escoamentos ideais,
pressão, rapidez e altura trocam energia entre si; perdas reais aparecem por viscosidade e turbulência.
\end{copiar}

\begin{dica}
Não confunda pressão com força. Pressão é força distribuída por área; dois êmbolos podem ter a
mesma pressão e forças diferentes.
\end{dica}

\begin{exercicios}
\nivel{1}{Conceitos}
\begin{questoes}
\item Calcule a densidade de $2\un{kg}$ ocupando $\dec{0,001}\un{m^3}$.
\item Uma força de $100\un{N}$ atua em $\dec{0,50}\un{m^2}$. Ache a pressão.
\item Em um mesmo líquido, qual ponto tem maior pressão: a $1\un{m}$ ou a $3\un{m}$ de profundidade?
\item Enuncie o princípio de Arquimedes.
\item O que ocorre com a rapidez quando a seção de um tubo diminui?
\nivel{2}{Aplicação}
\item Na água ($\rho=1000\un{kg/m^3}$), calcule $\Delta p$ a $2\un{m}$ de profundidade, com $g=10\un{m/s^2}$.
\item Um corpo desloca $\dec{0,030}\un{m^3}$ de água. Ache o empuxo.
\item Uma prensa tem $A_2=20A_1$. Qual força aparece no êmbolo maior se $F_1=150\un{N}$?
\item Água escoa em tubo de área $4\un{cm^2}$ a $2\un{m/s}$. Em área $2\un{cm^2}$, qual é a rapidez?
\nivel{3}{Síntese}
\item Por que uma barragem costuma ser mais espessa perto da base?
\item Um objeto de densidade média menor que a da água flutua. Explique usando peso e empuxo.
\item Uma mangueira estreitada lança água mais rapidamente. Relacione isso à continuidade.
\end{questoes}
\gabarito{1) $2000\un{kg/m^3}$; 2) $200\un{Pa}$; 3) a $3\un{m}$;
5) aumenta; 6) $2\times10^4\un{Pa}$; 7) $300\un{N}$; 8) $3000\un{N}$;
9) $4\un{m/s}$.}
\end{exercicios}


\end{document}
